\chapter{Problem Identification And Objectives}
\section{Problem Statement}

Although large language models are capable of generating high-quality code snippets, most current tools treat code generation as a single-step transformation from a prompt to source files. This approach lacks a systematic process for requirement analysis, architecture design, iterative refinement, and validation, leading to inconsistent project structures and limited support for maintenance. There is no explicit notion of roles, responsibilities, or quality gates, and generated projects often fail to meet production standards in terms of testing, modularity, and documentation.

The problem addressed in this project is to design and implement an \textbf{agentic code generation system} that can transform natural language specifications into complete, production-ready codebases using a structured, multi-agent pipeline. The system should coordinate specialized agents that emulate the roles of a human software team and should expose clear APIs and user interfaces to create, inspect, and iteratively improve projects generated by the pipeline.

\section{Objectives}

The primary objective of this project is to develop \textbf{MetaGPT}, a production-ready, SOP-driven agentic system for code generation. The specific objectives are:
\begin{itemize}
    \item To design a multi-agent pipeline consisting of Manager, Architect, Engineer, and QA agents that operate on structured inputs and outputs.
    \item To implement a FastAPI backend that exposes RESTful endpoints for project creation, pipeline execution, file management, and chat-based iteration.
    \item To integrate Google Gemini via LangChain and LangGraph for orchestrating the agent workflow and maintaining pipeline state.
    \item To develop a Next.js frontend that provides a dark-themed user interface for viewing agent outputs, browsing generated files, monitoring pipeline execution, and interacting with projects.
    \item To provide an isolated \textbf{live preview} capability for generated projects using an external sandbox service with non-blocking creation and status polling.
    \item To persist project metadata, pipeline status, and agent outputs in a \textbf{MongoDB} database to enable recovery, multi-user access control, and consistent iteration over time.
    \item To evaluate the system in terms of usability, consistency of generated project structure, and support for iterative refinement compared to one-shot LLM-based tools.
\end{itemize}

\section{Scope and Applications}

The scope of MetaGPT primarily targets full-stack web application generation, where the system can detect and scaffold common frameworks such as React and Next.js on the frontend and FastAPI-based services on the backend. The architecture, however, is intentionally generic and can be extended to support other project templates, languages, and deployment targets by updating the agent SOPs and project templates.

In the current implementation, the scope includes operational features required for practical usage: persistent storage of project state in MongoDB and optional live preview of generated projects through an isolated sandbox environment. These capabilities support longer workflows, interactive debugging, and reliable project access across application restarts.

The system can be applied in several scenarios:
\begin{itemize}
    \item Rapid prototyping of new product ideas by generating initial codebases from high-level prompts.
    \item Educational settings, where students can observe how requirements are translated into architecture, implementation, and tests by different agents.
    \item Internal tooling for organizations that wish to standardize how LLMs are used to generate and refactor code according to established conventions.
    \item As a foundation for future research on agentic software engineering, including automatic refactoring, performance optimization, and continuous integration with human feedback.
\end{itemize}

\section{Feasibility Analysis}

The feasibility of MetaGPT can be analyzed across technical, economic, and operational dimensions:
\begin{itemize}
    \item \textbf{Technical Feasibility}: The backend is implemented in Python 3.11 using FastAPI, which is well-suited for building high-performance, asynchronous APIs. Integration with Google Gemini is achieved through LangChain and LangGraph, both of which provide mature abstractions for LLM orchestration and agent workflows. The frontend is built with Next.js 14, TypeScript, and Tailwind CSS, technologies that are widely adopted and supported by the developer community.
    \item \textbf{Economic Feasibility}: The system is designed to run on commodity hardware with no requirements for proprietary infrastructure beyond the LLM API key. The use of open-source frameworks (FastAPI, Next.js, Tailwind CSS, and others) reduces licensing costs. Depending on usage patterns, LLM inference costs can be controlled by batching requests, caching intermediate artifacts, and tuning model parameters.
    \item \textbf{Operational Feasibility}: The project structure separates backend and frontend concerns, making it straightforward to deploy and maintain the system. Clear environment configuration files, documentation, and scripts for testing, linting, and type checking support day-to-day operations. The RESTful API and web interface make the system accessible to both developers and non-technical stakeholders.
\end{itemize}

\section{Project Plan}

The project was executed using a phased, iterative plan that covered analysis, design, implementation, testing, and documentation. Each phase had clearly defined deliverables, timelines, and responsibilities.

\subsection*{Phase I -- Requirement Analysis and High-Level Design (Week 1--2)}
\begin{itemize}
    \item Conducted a survey of existing LLM-based code generation tools and agentic frameworks.
    \item Refined the problem statement, objectives, scope, and feasibility based on literature review and initial experiments.
    \item Identified primary user personas (developer, student, researcher) and their interaction flows with the system.
    \item Drafted the high-level architecture, including the separation into backend, agent pipeline, storage, and frontend subsystems.
    \item Defined initial success criteria (e.g., ability to generate a multi-file React/FastAPI project from a single prompt).
\end{itemize}

\subsection*{Phase II -- Backend and Agent Pipeline Design (Week 3--4)}
\begin{itemize}
    \item Designed the project structure under \texttt{backend/\allowbreak{}app} including \texttt{api}, \texttt{agents}, \texttt{graph}, \texttt{services}, \texttt{storage}, \texttt{schemas}, and \texttt{rag}.
    \item Specified the responsibilities and data contracts for the Manager, Architect, Engineer, and QA agents.
    \item Defined the LangGraph pipeline state in \texttt{graph/\allowbreak{}\allowbreak{}state.py} and documented the transitions in \texttt{graph/\allowbreak{}\allowbreak{}orchestrator.py}.
    \item Designed REST API endpoints for projects, pipeline execution, files, chat, and RAG operations in the \texttt{api/\allowbreak{}\allowbreak{}endpoints} package.
    \item Finalized environment configuration parameters in \texttt{config.py} and designed .env templates.
\end{itemize}

\subsection*{Phase III -- Backend Implementation (Week 5--7)}
\begin{itemize}
    \item Implemented the FastAPI application entry point in \texttt{main.py} and wired the routers defined in \texttt{api/\allowbreak{}\allowbreak{}router.py}.
    \item Implemented agent classes in \texttt{agents/\allowbreak{}\allowbreak{}manager.py}, \texttt{agents/\allowbreak{}\allowbreak{}architect.py}, \texttt{agents/\allowbreak{}\allowbreak{}engineer.py}, and \texttt{agents/\allowbreak{}\allowbreak{}qa.py}, following SOPs from \texttt{sop/\allowbreak{}\allowbreak{}definitions.py}.
    \item Developed pipeline and chat services in \texttt{services/\allowbreak{}\allowbreak{}pipeline\_\allowbreak{}service.py} and \texttt{services/\allowbreak{}\allowbreak{}chat\_\allowbreak{}service.py}.
    \item Implemented storage backends for projects and files in \texttt{storage/\allowbreak{}\allowbreak{}project\_\allowbreak{}store.py} and \texttt{storage/\allowbreak{}\allowbreak{}file\_\allowbreak{}store.py}.
    \item Integrated Google Gemini configuration via LangChain in \texttt{llm/\allowbreak{}\allowbreak{}gemini.py}, including helper functions for structured output.
\end{itemize}

\subsection*{Phase IV -- Frontend Implementation (Week 8--10)}
\begin{itemize}
    \item Created the Next.js App Router structure under \texttt{frontend/\allowbreak{}\allowbreak{}src/\allowbreak{}app}, including landing, projects list, project workspace, and authentication pages.
    \item Implemented core UI components such as \texttt{AgentOutputs}, \texttt{ChatPanel}, \texttt{CodeViewer}, \texttt{ExecutionTimeline}, and \texttt{FileExplorer}.
    \item Set up global styling and theming with Tailwind CSS and dark theme defaults.
    \item Implemented API client utilities in \texttt{frontend/\allowbreak{}\allowbreak{}src/\allowbreak{}lib/\allowbreak{}api.ts} and application state management using stores in \texttt{frontend/\allowbreak{}\allowbreak{}src/\allowbreak{}lib/\allowbreak{}store.ts}.
    \item Connected frontend interactions (project creation, pipeline run, file viewing, chat) to backend endpoints.
\end{itemize}

\subsection*{Phase V -- Testing, Evaluation, and Refinement (Week 11--12)}
\begin{itemize}
    \item Wrote and executed unit tests for backend services, schemas, and utility functions using \texttt{pytest}.
    \item Performed integration tests covering project lifecycle, pipeline execution, and file retrieval operations.
    \item Ran linting and type-checking for both backend and frontend (\texttt{ruff}, \texttt{mypy}, ESLint, TypeScript).
    \item Evaluated the system using representative prompts and collected observations about code quality, structure, and usability.
    \item Refined agent prompts, SOP definitions, and UI feedback based on test results and manual inspection.
\end{itemize}

\subsection*{Phase VI -- Documentation and Deployment Readiness (Week 13)}
\begin{itemize}
    \item Prepared project documentation including top-level \texttt{README.md}, backend and frontend READMEs, and this report.
    \item Documented environment setup steps, configuration options, and development workflows.
    \item Verified that the project can be cloned, configured, and run from a clean environment using the documented instructions.
    \item Identified future enhancements and roadmap items for subsequent iterations.
\end{itemize}
